\newpage
\section{Manual de usuario}

Este manual é autocontido e está dirixido ás persoas que utilicen Telaria. Describe a
instalación e o uso de todas as funcionalidades. As capturas de pantalla insírense nos
espazos reservados ao efecto.

\subsection{Instalación e acceso}

Telaria é unha aplicación móbil para iOS e Android. Durante o desenvolvemento pode
executarse a través de Expo Go ou instalando unha compilación da aplicación. Ao abrila por
primeira vez, o usuario debe rexistrarse ou iniciar sesión.

\subsection{Rexistro e inicio de sesión}

Na pantalla inicial pódese crear unha conta nova indicando nome de usuario, correo
electrónico e contrasinal (con confirmación), ou iniciar sesión cunha conta existente. A
sesión mantense aberta entre usos.

\figcaptura{Pantallas de rexistro e inicio de sesión}{fig:b-login}

\subsection{Pantalla principal: as miñas viaxes}

Tras iniciar sesión, móstrase a lista de viaxes do usuario. Desde aquí pódese acceder a
unha viaxe, crear unha nova ou unirse a unha existente, así como abrir o menú lateral para
acceder a amigos, invitacións, perfil e axustes.

\figcaptura{Lista de viaxes e menú principal}{fig:b-main}

\subsection{Crear e unirse a viaxes}

Para crear unha viaxe abonda con indicar o seu nome; opcionalmente pódense convidar amigos
no momento da creación. Para unirse a unha viaxe existente pódese empregar unha invitación
recibida, enviar unha solicitude de unión ou escanear un código QR.

\figcaptura{Creación de viaxe e unión mediante código QR}{fig:b-crear}

\subsection{Xestión de membros}

Dentro dunha viaxe, calquera membro pode convidar outros usuarios e xestionar as
solicitudes de unión pendentes. Non existen roles: todos os membros teñen as mesmas
capacidades.

\figcaptura{Membros da viaxe e solicitudes pendentes}{fig:b-membros}

\subsection{Gastos e liquidacións}

No apartado de gastos rexístranse os pagamentos indicando importe, concepto, categoría,
quen pagou e entre quen se reparte. A aplicación calcula automaticamente os balances (quen
debe a quen) e propón as liquidacións que minimizan o número de transferencias. As débedas
saldadas rexístranse no histórico de liquidacións.

\figcaptura{Lista de gastos, balances e liquidacións}{fig:b-gastos}

\subsection{Eventos e calendario}

O calendario da viaxe permite crear eventos con nome, data e hora de inicio, duración e
unha localización opcional. Os eventos amósanse organizados por data.

\figcaptura{Calendario e creación de eventos}{fig:b-eventos}

\subsection{Documentos compartidos}

Os membros poden subir documentos relevantes para a viaxe (reservas, billetes, etc.) e
consultalos ou eliminalos. Os documentos de imaxe amosan unha miniatura de previsualización.

\figcaptura{Listaxe de documentos compartidos}{fig:b-documentos}

\subsection{Chat de grupo}

Cada viaxe dispón dun chat común no que os membros intercambian mensaxes en tempo real.

\figcaptura{Chat de grupo da viaxe}{fig:b-chat}

\subsection{Asistente de intelixencia artificial}

Cada viaxe inclúe un asistente conversacional que responde preguntas tendo en conta o
contexto da viaxe (eventos e gastos). A conversa garda un historial persistente.

\figcaptura{Asistente de IA da viaxe}{fig:b-ia}

\subsection{Amizades}

A sección de amigos permite enviar e xestionar solicitudes de amizade (por identificador
ou código QR), consultar a lista de amigos e eliminar amizades. Os amigos poden convidarse
ás viaxes de forma máis áxil.

\figcaptura{Lista de amigos e solicitudes de amizade}{fig:b-amigos}

\subsection{Perfil e axustes}

No perfil pódese consultar e actualizar o nome de usuario e o correo electrónico (o cambio
de correo require o contrasinal), cambiar o contrasinal e establecer un avatar. Nos axustes
configúrase o idioma (galego, castelán ou inglés), o tema (claro ou escuro) e o aforro de
datos.

\figcaptura{Perfil de usuario e pantalla de axustes}{fig:b-perfil}

\subsection{Eliminación da conta}

Desde os axustes, o usuario pode eliminar a súa conta. A operación require a confirmación
do contrasinal e elimina os datos persoais asociados.

\figcaptura{Confirmación de eliminación de conta}{fig:b-eliminar}

\subsection{Mensaxes de erro comúns}

\begin{itemize}
  \item \textbf{Credenciais incorrectas}: o correo ou o contrasinal non son válidos no
        inicio de sesión.
  \item \textbf{Correo xa rexistrado}: existe xa unha conta con ese correo ao rexistrarse.
  \item \textbf{Acción xa realizada}: por exemplo, o usuario xa é membro da viaxe ou xa foi
        convidado.
  \item \textbf{Sen conexión}: a aplicación require conexión á rede; comprobe a
        conectividade e que o servidor estea accesible.
\end{itemize}
